\section{结论}
\label{sec:conclusion}

当监测关注的是传播的典型水平或下一时期的取值时，基于病例数的再生数估计，其精度随病例量增长而递减地提高。超过平衡规模后，同一时期内共有的变异占主导，病例数每再加倍一次，至多只能消除剩余方差的四分之一。在美国流感住院数据中，按更新型比值的尺度，这一规模约为每州每周 37 例住院。以 3.6 天的代际间隔和相互独立的周冲击为条件，经模拟校准后隐含的环境方差约为 0.04，共有变异中相当一部分为同一周内各州所共有。对每期报告病例已超过平衡规模的监测系统，增加观测时期、增加变异相互独立的区域，比增加每期病例更能提高精度。关于报告的证据范围较窄：强制报告下，报告完整度对截距的贡献很小；在观察性比较中，较高的机构参与度与波动较小的增长统计量相伴。表观统计量的经验截距应与联系其与传播的校准一并报告。
